\section{ANEXOS}

Los anexos acompañan a esta memoria en la carpeta Anexos de la entrega, cada uno con su fuente
editable y su PDF. La Tabla~\ref{tab:anexos} los lista con su contenido y su estado. Los
documentos técnicos de soporte (VFP, SRS, SPMP, WBS, SAD, SDD, documento maestro de diseño,
informes ATAM y QA, informe de pruebas, que está en la carpeta QA, Dossier de OE1 y diagramas
de OE2) están en sus carpetas de la entrega y se citan en el texto cuando se usan.

\begin{table}[htbp]
\centering
\caption{Índice de anexos de la entrega}
\label{tab:anexos}
\small
\renewcommand{\arraystretch}{1.2}
\begin{tabular}{@{}p{4.2cm}p{7.4cm}p{2.4cm}@{}}
\toprule
\textbf{Anexo} & \textbf{Contenido} & \textbf{Estado} \\
\midrule
Registros de decisión (ADR) & Treinta y ocho registros de decisión de arquitectura (0001 a 0038),
  en formato Nygard, reunidos en un solo documento & Entregado \\
\addlinespace
Actas de aprobación & Actas de los hitos OE1 y OE2 & Entregadas; sin firmas \\
\addlinespace
Auditoría de autorización & Autorización de los servicios anteriores y del sistema construido,
  con la situación de cada hallazgo & Entregado \\
\addlinespace
Anexo M, datos personales & Registro de tratamiento de datos personales del sistema
  construido, con la clasificación de los datos de menores y su base legal (Ley 1581 de 2012) &
  Entregado \\
\addlinespace
Diccionario de datos & Diccionario de datos y esquema objetivo & Entregado \\
\addlinespace
Empatía y entrevistas & Guiones de entrevista y mapas de empatía por actor & Entregado \\
\addlinespace
Matriz de evidencia & Matriz de objetivos específicos y evidencia & Entregado \\
\addlinespace
Priorización MoSCoW & Priorización de los requisitos del SRS v6.1, con su estado en el sistema
  construido & Entregado \\
\addlinespace
CRISP-DM de recuperación & Recuperación ordenada por las seis fases de CRISP-DM, con la
  evolución desde los prototipos y el estado, las pruebas y la justificación de cada
  estrategia & Entregado \\
\addlinespace
Diagramas de clases & Clases de los servicios y de la biblioteca común & Entregado \\
\midrule
Previstos en la WBS & Protocolo ético y consentimientos, instrumentos de validación con
  usuarios, kits e hipervolumen del frente de Pareto y demostración del incremento con su
  retrospectiva. La cobertura se reporta en el informe de pruebas & Se incorporan con la entrega final \\
\bottomrule
\end{tabular}
\end{table}

El protocolo ético, los consentimientos y los instrumentos de validación corresponden al OE5
(Sección~\ref{sec:oe5})\pendienteEvidencia{anexo de protocolo etico, consentimiento e instrumentos del OE5}. Los diagramas de contexto y de contenedores
de la alternativa elegida son las Figuras~\ref{fig:c4-contexto} y~\ref{fig:c4-contenedores};
los de las tres alternativas evaluadas están en el informe ATAM y en los diagramas de OE2, y
los de componentes y de despliegue, en el SAD.

\subsection{Tabla de constantes numéricas}
\label{anx:constantes}

La Tabla~\ref{tab:constantes} reúne cada constante numérica del sistema con su valor, su módulo, su origen y su estado, según se explica en la Sección~\ref{sec:constantes}.

{\footnotesize
\setlength{\tabcolsep}{3pt}
\renewcommand{\arraystretch}{1.1}
\begin{longtable}{@{}p{3.1cm}p{2.3cm}p{2.5cm}p{4.7cm}p{1.6cm}@{}}
\caption{Constantes numéricas del sistema, con su origen y su estado.}\label{tab:constantes}\\
\toprule
Constante & Valor & Módulo & Origen & Estado \\
\midrule
\endfirsthead
\toprule
Constante & Valor & Módulo & Origen & Estado \\
\midrule
\endhead
\bottomrule
\endfoot
\multicolumn{5}{@{}l}{\emph{Percepción}}\\
Tope de páginas & 40 & percepción & Experimento: el PDF más largo del corpus tiene 39 páginas & Vigente \\
Temperatura del modelo & 0 & percepción & Decodificación voraz, para reducir la variabilidad entre lecturas & Vigente \\
Tiempo límite por llamada & 40 s & percepción & Una página típica tarda unos 5 s; 40 s indica una llamada colgada & Vigente \\
Reintentos y espera & 1 reintento, 3 s & percepción & Solo ante fallas de transporte, para no ocultar respuestas mal formadas & Vigente \\
\midrule
\multicolumn{5}{@{}l}{\emph{Normalización y resolución}}\\
Dimensión del vector & 1.024 & catálogo & Fija la columna de la base; se pide explícitamente al modelo de \emph{embeddings} & Vigente \\
Lote de \emph{embeddings} & 100 textos & catálogo & Una petición por lote, para evitar cuerpos de petición excesivos & Vigente \\
Umbral de reconocimiento & 0,5 & catálogo & Corte por debajo del cual el renglón no corresponde a ningún concepto; fijado al corregir el caso «chaquiras» & Vigente \\
Aporte del coseno $\gamma$ & 0,15 & catálogo & Ajuste sobre casos de error; pequeño para que el coseno desempate y no decida & En ciclo de calibración \\
Umbral de rescate $\theta$ & 0,86 & catálogo & Ajuste sobre casos de error; separa un sinónimo de un parecido de tema & En ciclo de calibración \\
Piso del rescate $\phi$ & 0,55 & catálogo & Apenas sobre el corte de 0,5, para que un rescate entre a revisión & En ciclo de calibración \\
Umbral de autoaceptación $\tau$ & 0,75 & inteligencia & Requisito PE-04 del SRS; valor inicial anterior a la calibración (0,85 y 0,630 en versiones previas); ver Tabla~\ref{tab:cal-optim} & En ciclo de calibración \\
Tolerancia de medidas & 0,25 y rango 0,85 a 1,35 & catálogo & Ajuste con las auditorías de los kits (Tabla~\ref{tab:kits}) & Vigente \\
Rango de precio en tiendas & \$1.000 a \$400.000 & catálogo & Descarta accesorios sueltos y lotes institucionales & Vigente \\
\midrule
\multicolumn{5}{@{}l}{\emph{Optimización}}\\
Límite de enumeración & 4.096 ($2^{12}$) & inteligencia & Por debajo de ese tamaño del espacio el frente exacto se calcula en milisegundos; por encima se usa NSGA-II & Vigente \\
Población & $\min(200,\max(40,12n))$, par & inteligencia & Escala con las $n$ líneas decidibles; el par evita desperdiciar un individuo & Vigente \\
Generaciones & $\min(120,\max(30,20\log_{10}|\mathcal{E}|))$ & inteligencia & Escala con el logaritmo del espacio $|\mathcal{E}|$ de búsqueda & Vigente \\
Cruce SBX & $p_c=0{,}9$, $\eta_c=15$ & inteligencia & Valores estándar de \citet{deb_simulated_1995} & Vigente \\
Mutación polinomial & $p=1$, $\eta_m=20$ & inteligencia & Valores estándar de \citet{deb_combined_1996} & Vigente \\
Semilla & 0 & inteligencia & Reproducibilidad del frente entre ejecuciones & Vigente \\
Pesos por omisión (costo, marca, calidad, personaje) & 0,35; 0,30; 0,25; 0,10 & inteligencia & Encuesta, pregunta P21 (n=60), normalizada; ver Tabla~\ref{tab:cal-optim} & En ciclo de calibración \\
Sobreprecio para Premium & 10\,\% & inteligencia & Encuesta: 42 de 60 pagan como máximo un 5\,\% más; el 10\,\% marca a quien acepta pagar más & Vigente \\
Opciones por renglón & 3 & catálogo & Pocas opciones diferenciadas facilitan la decisión \citep{iyengar_when_2000} & Vigente \\
\midrule
\multicolumn{5}{@{}l}{\emph{Requisitos de desempeño y operación}}\\
Tiempo de generación del kit & 8 s & SRS (DP-03) & Requisito de desempeño; se contrasta en la Sección~\ref{sec:optimizacion} & Vigente \\
Ventana de revisión & 4 h en temporada & SRS (ML-09) & Requisito: el administrador revisa cada lista nueva antes de publicarla & Vigente \\
Disyuntor & 5 fallas, 30 s & biblioteca común & Patrón de resiliencia; 30 s es el escenario ESC-R2 del ATAM & Vigente \\
Caché de lecturas & 30 días & inteligencia & Evita pagar de nuevo la percepción de un mismo documento & Vigente \\
Caché de vectores & 24 h & catálogo & Evita recalcular el vector de un renglón repetido & Vigente \\
\end{longtable}
}

\subsection{Glosario}
\label{anx:glosario}

\begin{description}[leftmargin=2.5cm,style=nextline]
\item[MAUT] Teoría de utilidad multiatributo: método que resume varios criterios en un solo valor de utilidad mediante pesos que suman uno.
\item[NSGA-II] Algoritmo genético de ordenamiento no dominado que aproxima el frente de Pareto de un problema multiobjetivo \citep{deb_fast_2002}.
\item[Frente de Pareto] Conjunto de soluciones no dominadas, es decir, aquellas en las que ninguna otra mejora un objetivo sin empeorar otro.
\item[Hipervolumen] Volumen del espacio de objetivos que un frente domina respecto de un punto de referencia; mide la calidad de un conjunto de soluciones \citep{zitzler_multiobjective_1999}.
\item[NL2C] Traducción de una preferencia en lenguaje natural a restricciones estructuradas que el optimizador entiende.
\item[$\tau$] Umbral de autoaceptación: puntaje por encima del cual un renglón se da por resuelto sin revisión individual.
\item[$\delta$] Margen mínimo entre el mejor y el segundo candidato que usó una versión anterior del umbral.
\item[$\kappa$] Coeficiente kappa de Cohen: acuerdo entre dos clasificadores corregido por el azar \citep{cohen_coefficient_1960}.
\item[SUS] Escala de usabilidad del sistema, de 0 a 100 \citep{brooke_sus_1996}.
\item[NPS] Índice de recomendación neta, de $-100$ a $100$ \citep{reichheld_one_2003}.
\item[RTLX] Versión sin ponderar del índice de carga de trabajo NASA-TLX, de 0 a 100.
\item[CRISP-DM] Proceso de seis etapas para proyectos de datos \citep{chapman_crisp-dm_2000,wirth_crisp-dm_2000}.
\item[ATAM] Método de evaluación de arquitecturas por compromisos entre atributos de calidad \citep{bass_software_2021}.
\item[ADR] Registro de decisión de arquitectura.
\item[Concepto canónico] Clase de productos equivalentes del catálogo a la que se asigna un renglón de la lista.
\end{description}
